% fig_lec09_sphere_holonomy.tex — Path-dependent parallel transport on S²
% Data: generated by scripts/sim_lec09_lec10.jl
%
% Initial vector w₀ = (1,0,0) at the north pole p = (0,0,1).
% Two paths from p to q = (1,0,0):
%   γ₁ : meridian φ = 0                   p → q
%   γ₂ : meridian φ = π/2 then equator    p → (0,1,0) → q
% Parallel transport along each path is integrated as the projection ODE
%   dw/ds = −(w · γ̇) γ(s)        on the unit sphere.
% Final vectors at q:  w_{γ₁} = (0,0,−1),  w_{γ₂} = (0,−1,0)  — they
% differ by π/2 (= the area of the spherical triangle they bound).

\begin{center}
\begin{tikzpicture}
  \begin{axis}[
      width=10.0cm, height=8.4cm,
      view={-32}{22},
      xmin=-1.15, xmax=1.15, ymin=-1.15, ymax=1.15, zmin=-1.15, zmax=1.15,
      axis equal image,
      axis lines=none,
      xtick=\empty, ytick=\empty, ztick=\empty,
      clip=false,
  ]
    % Sphere (faded)
    \addplot3[surf, opacity=0.28, shader=interp,
              colormap name=grsurface, mesh/rows=50, forget plot]
      table {data/lec09_sphere.dat};

    % Path γ₁ (meridian φ=0)
    \addplot3[munsell, line width=1.3pt, no markers, forget plot]
      table {data/lec09_path1_curve.dat};

    % Path γ₂ (meridian φ=π/2 + equator)
    \addplot3[cgblue!75!black, line width=1.3pt, no markers, forget plot]
      table {data/lec09_path2_curve.dat};

    % Sample transported vectors along γ₁
    \addplot3[
      quiver={u=\thisrowno{3}, v=\thisrowno{4}, w=\thisrowno{5}},
      -{Stealth[length=4.5pt, width=3pt]},
      munsell!70!black, line width=0.85pt, no markers, forget plot,
    ] table[x index=0, y index=1, z index=2]
      {data/lec09_path1_vectors.dat};

    % Sample transported vectors along γ₂
    \addplot3[
      quiver={u=\thisrowno{3}, v=\thisrowno{4}, w=\thisrowno{5}},
      -{Stealth[length=4.5pt, width=3pt]},
      cgblue!75!black, line width=0.85pt, no markers, forget plot,
    ] table[x index=0, y index=1, z index=2]
      {data/lec09_path2_vectors.dat};

    % Initial vector w₀ at p (north pole)
    \addplot3[
      quiver={u=\thisrowno{3}, v=\thisrowno{4}, w=\thisrowno{5}},
      -{Stealth[length=5.5pt, width=4pt]},
      banana!75!black, line width=1.2pt, no markers, forget plot,
    ] table[x index=0, y index=1, z index=2]
      {data/lec09_v_initial.dat};

    % Final w along γ₁ at q
    \addplot3[
      quiver={u=\thisrowno{3}, v=\thisrowno{4}, w=\thisrowno{5}},
      -{Stealth[length=5.5pt, width=4pt]},
      munsell!75!black, line width=1.4pt, no markers, forget plot,
    ] table[x index=0, y index=1, z index=2]
      {data/lec09_v_path1_at_q.dat};

    % Final w along γ₂ at q
    \addplot3[
      quiver={u=\thisrowno{3}, v=\thisrowno{4}, w=\thisrowno{5}},
      -{Stealth[length=5.5pt, width=4pt]},
      cgblue!85!black, line width=1.4pt, no markers, forget plot,
    ] table[x index=0, y index=1, z index=2]
      {data/lec09_v_path2_at_q.dat};

    % Marker points: p, q, and the (0,1,0) corner
    \addplot3[only marks, mark=*, mark size=2.0pt, banana!90!black, forget plot]
      table {data/lec09_marker_points.dat};

    % Labels
    \node[font=\sf\small, text=spacecadet]
      at (axis cs:-1.10,-0.95,-0.50) {$S^2$};
    \node[font=\small, text=banana!60!black, anchor=south east]
      at (axis cs:0.0,0.0,1.04) {$p$};
    \node[font=\small, text=banana!60!black, anchor=south west]
      at (axis cs:1.05,0.0,0.05) {$q$};
    \node[font=\scriptsize, text=banana!50!black, anchor=west]
      at (axis cs:0.30,0.0,1.02) {$w_0$};
    \node[font=\scriptsize, text=munsell, anchor=west]
      at (axis cs:0.95,0.0,0.50) {$\gamma_1$};
    \node[font=\scriptsize, text=cgblue!75!black, anchor=south east]
      at (axis cs:0.0,0.85,0.55) {$\gamma_2$};
    \node[font=\scriptsize, text=munsell!70!black, anchor=west]
      at (axis cs:1.10,0.05,-0.35) {$w_{\gamma_1}\!(q)$};
    \node[font=\scriptsize, text=cgblue!85!black, anchor=west]
      at (axis cs:1.10,-0.40,-0.02) {$w_{\gamma_2}\!(q)$};
  \end{axis}
\end{tikzpicture}
\end{center}
